\documentclass[../notes.tex]{subfiles}
\begin{document}
    
    \subsection{Estructura}
    
        Beamer es una clase de documento para realizar presentaciones.

            \begin{verbatim}
\documentclass{beamer}
            \end{verbatim}

        La estructura básica de un beamer corresponderán a slides, de forma que denotaremos \texttt{frame} a cada diapositiva de la presentación.
        
        \begin{verbatim}
\begin{frame}{Titulo de la diapositiva}
    Aqui va el texto de la diapositiva.
\end{frame}
        \end{verbatim}

        En general, siguen las mismas reglas vistas anteriormente. Beamer entrega más herramientas para utilizar.

    
        Se pueden también crear título e índices para las presentaciones.
        \begin{verbatim}
\begin{frame}
    \maketitle
\end{frame}

\begin{frame}
    \tableofcontents
\end{frame}
        \end{verbatim}
        
%        \begin{block}{Nota}
        	Existen más opciones para agregar al título que los vistos anteriormente. En este caso, también podemos agregar:
%        \end{block}
        \begin{verbatim}
\subtitle{Aquí va un texto corto}
\logo{\includegraphics{utfsm}}
        \end{verbatim}
        
        También se puede crear un índice donde destaque la sección y sub-sección actual.
        \begin{verbatim}
\begin{frame}
    \tableofcontents[currentsection]
\end{frame}
\begin{frame}
    \tableofcontents[currentsection,currentsubsection]
\end{frame}
        \end{verbatim}

    \subsection{Múltiples columnas en diapositivas}

        Se puede utilizar múltiples columnas con el ambiente \texttt{columns}
        \begin{verbatim}
\begin{columns}
    \column{0.4 \linewidth}
        Este es el texto de la columna 1.
            
        Puede contener varios párrafos.
    \column{0.6 \linewidth}
        Este es el texto de la columna 2.
            
        Esta es una columna más ancha que la otra.
\end{columns}
        \end{verbatim}
    Output:
    
        %\begin{columns}
        %    \column{0.4 \linewidth}
        %        Este es el texto de la columna 1.
        %
        %        Puede contener varios párrafos.
        %    \column{0.6 \linewidth}
        %        Este es el texto de la columna 2.
        %
        %        Esta es una columna más ancha que la otra.
        %\end{columns}
        
        %\begin{block}{Nota}
            Tiene sentido establecer el largo de la columna de forma que dependa del largo de la linea. En este caso, se estableció que la primera columna posea el 40\% del largo de la linea y la segunda columna el 60\%.\\[\baselineskip]
            
            No necesariamente debe cubrir toda la linea.
        %\end{block}
        
        \begin{verbatim}
    Este texto se encuentra arriba.
        \vspace{1em}
\begin{columns}
    \column{0.4 \linewidth}
        Este es el texto de la columna 1.
            
        Puede contener varios párrafos.
    \column{0.4 \linewidth}
        Este es el texto de la columna 2.
            
        Esta es una columna más ancha que la otra.
\end{columns}
        \vspace{1em}
    Este texto se encuentra abajo.
        \end{verbatim}
        
    Este texto se encuentra arriba de las columnas.
    \vspace{1em}
%\begin{columns}
%    \column{0.4 \linewidth}
        Este es el texto de la columna 1.
            
        Puede contener varios párrafos.
%    \column{0.4 \linewidth}
        Este es el texto de la columna 2.
            
        Esta es una columna más ancha que la otra.
%\end{columns}
    \vspace{1em}
    Este texto se encuentra abajo de las columnas.

    \subsection{Pausas}
    
    	Una diapositiva se puede separar en varias a medida que se agregan pausas: hacer aparecer una linea nueva en la slide cada vez que se avanza en la presentación.
        \begin{verbatim}
\begin{frame}
    Este texto aparecerá en una primera instancia.
    
    \pause
    
    Este texto aparecerá en una segunda instancia.
\end{frame}
        \end{verbatim}

    \subsection{Capas especificas}
    
    	Además de las pausas, se pueden especificar las slides en las que aparecerán, utilizando \texttt{<>}.
        \begin{verbatim}
\begin{frame}
    \begin{itemize}
        \item<2> Esta solo aparecerá en la slide 2
        \item<1-4> Esta aparecerá desde la slides 1 
        hasta la 4, incluyéndolas.
        \item<3-> Esta aparecerá desde la slide 3 en adelante.
        \item<1,3-4,6-> Esta aparecerá en la slide 1, 
        desde la 3 hasta la 4 y desde la 6 en adelante.
    \end{itemize}
\end{frame}
        \end{verbatim}
    
    \subsection{Bloques}
        \begin{verbatim}
\begin{frame}
    \begin{block}{Titulo del bloque}
        Este texto se encuentra dentro del bloque.
    \end{block}
    \begin{alertblock}{Bloque de alerta}
        Este texto se encuentra dentro de un bloque de alerta.
    \end{block}
\end{frame}
        \end{verbatim}
        %\begin{block}{Titulo del bloque}
        %    Este texto se encuentra dentro del bloque.
        %\end{}
        %\begin{alertblock}{Bloque de alerta}
        %    Este texto se encuentra dentro de un bloque de alerta.
        %\end{alertblock}
    
    \subsection{Temas}
        Para seleccionar un tema, se deben agregar los comandos \texttt{usetheme} y \texttt{usecolortheme} en el preámbulo.
        \begin{verbatim}
        	\usetheme{Frankfurt}
        \end{verbatim}
        
        %\begin{block}{Nota}
        	Los temas de beamer poseen nombres de ciudades; los temas de color poseen nombres de animales.
        %\end{block}

        Existen muchos temas disponibles para beamer. Se puede encontrar un listado de estas en
            \begin{itemize}
            	\item \href{https://www.overleaf.com/learn/latex/Beamer%23Creating_a_table_of_contents#Reference_guide}{Overleaf: guía de referencia}
            	\item \href{https://latex-beamer.com/tutorials/beamer-themes/}{Lista completa de Temas de Beamer}
            \end{itemize}
\end{document}
